% ---
% school: "college"
% level: "8eme"
% section: null
% chapter: "العمليات والحساب على الأعداد الكسرية"
% source: "Axiomaths"
% title: "قيمة عبارة والمقارنة"
% country: "Tunisie"
% year: 2026
% professor: "Axiomaths"
% difficulty: "Moyen"
% points: null
% tags: ["fractions", "expressions-litterales", "valeur-numerique", "comparaison"]
% ---
% Exercice 12/18 de la série 8ème « العمليات والحساب على الأعداد الكسرية » (statement + \axsol + solution). Compiler avec XeLaTeX + axiomath.sty (+ axiomaths-logo.png).
\documentclass[12pt,a4paper]{article}
\usepackage[level=8b,color=15803D]{axiomath}
\axsetup{title={العمليات والحساب على الأعداد الكسرية}, corrige}

\begin{document}

\begin{corr}[title={قيمة عبارة والمقارنة}, stars=3, source={Axiomaths}]
نعتبر العبارتين $E=2x+y-1$ و $F=x-3y$ حيث $x$ و $y$ عددان كسريان.
\begin{enumerate}
  \item احسب $E$ و $F$ في كل حالة :
  \begin{enumerate}
    \item $x=\dfrac{2}{3}$ و $y=-\dfrac{1}{2}$ ؛
    \item $x=\dfrac{3}{4}$ و $y=\dfrac{1}{2}$.
  \end{enumerate}
  \item قارن $E$ و $F$ في كل حالة من الحالتين.
  \item بين أن $E-F=x+4y-1$ مهما كان $x$ و $y$.
  \item تحقق من هذه المساواة في الحالتين (أ) و (ب). هل يمكن القول إن $E<F$ دائما ؟ علّل.
\end{enumerate}
\axsol
\begin{enumerate}[itemsep=14pt]
  \item نعوّض $x$ و $y$ بقيمتيهما ثم نوحّد المقامات.
  \begin{enumerate}
    \item $x=\dfrac{2}{3}$ و $y=-\dfrac{1}{2}$ :
    \[E=2\times\dfrac{2}{3}+\left(-\dfrac{1}{2}\right)-1=\dfrac{4}{3}-\dfrac{1}{2}-1=\dfrac{8}{6}-\dfrac{3}{6}-\dfrac{6}{6}=\dfrac{-1}{6}\]
    \[F=\dfrac{2}{3}-3\times\left(-\dfrac{1}{2}\right)=\dfrac{2}{3}+\dfrac{3}{2}=\dfrac{4}{6}+\dfrac{9}{6}=\dfrac{13}{6}\]
    \item $x=\dfrac{3}{4}$ و $y=\dfrac{1}{2}$ :
    \[E=2\times\dfrac{3}{4}+\dfrac{1}{2}-1=\dfrac{3}{2}+\dfrac{1}{2}-1=2-1=1\]
    \[F=\dfrac{3}{4}-3\times\dfrac{1}{2}=\dfrac{3}{4}-\dfrac{6}{4}=-\dfrac{3}{4}\]
  \end{enumerate}
  \item \begin{prop}[مقارنة عددين نسبيين]
  كل عدد سالب أصغر من كل عدد موجب.
  \end{prop}
  \begin{enumerate}
    \item $E=-\dfrac{1}{6}<0$ و $F=\dfrac{13}{6}>0$ ، إذن \ib{$E<F$}.
    \item $E=1>0$ و $F=-\dfrac{3}{4}<0$ ، إذن \ib{$E>F$}.
  \end{enumerate}
  \item عند حذف الأقواس المسبوقة بعلامة $(-)$ تتغير علامات ما بداخلها :
  \[E-F=(2x+y-1)-(x-3y)=2x+y-1-x+3y\]
  نجمع الحدود المتشابهة :
  \[\ib{$E-F=x+4y-1$}\]
  \item \begin{enumerate}
    \item $x+4y-1=\dfrac{2}{3}-2-1=\dfrac{2}{3}-3=-\dfrac{7}{3}$ ؛ و $E-F=-\dfrac{1}{6}-\dfrac{13}{6}=-\dfrac{14}{6}=-\dfrac{7}{3}$. \checkmark
    \item $x+4y-1=\dfrac{3}{4}+2-1=\dfrac{7}{4}$ ؛ و $E-F=1-\left(-\dfrac{3}{4}\right)=\dfrac{7}{4}$. \checkmark
  \end{enumerate}
  في الحالة (أ) $E-F<0$ فـ $E<F$ ، وفي الحالة (ب) $E-F>0$ فـ $E>F$. إذن \ib{لا} : الحالة (ب) مثال مضاد، فإشارة $E-F$ تتوقف على قيمتي $x$ و $y$.
\end{enumerate}
\textbf{الفكرة :} لمقارنة عبارتين حرفيتين نشكّل فرقهما ونبسّطه ؛ ثم نناقش إشارته حسب قيم المتغيرين.
\end{corr}

\end{document}
